\section{Reproducibility}\label{sec:repro}


The computation can be replicated at three levels: by verifying the certificate with the supplied
verifiers (seconds to minutes), by regenerating the certificate and all checks from the sources (under an
hour), or by writing a verifier from this paper alone. This section describes the supplementary archive
and gives the commands and the expected outputs; the contents of the archive, the complete outputs and
the running times that we measured are given in \cref{sec:supp-repro}.

\subsection{The archive}\label{ssec:artifact}

The supplementary archive is one file, \fn{ttt6-artifact.tar.gz} (\ArchiveSize; $\ArchiveFiles$ files,
\ArchiveUnpacked{} unpacked), produced by
\fn{scripts/make_artifact.sh}. It contains sources and results but no binaries. Its file
\fn{MANIFEST.sha256} lists the SHA-256 checksum of every other file. The contents are listed in
\cref{tab:artifact} in \cref{ssec:supp-archive}. Every file that is named in this paper is contained in
the archive, under the path given in the paper. The documents \fn{SPEC.md}, \fn{PROOF.md}, \fn{REPORT.md} and \fn{REVIEW.md} and the comments in the
sources predate this paper and call Maker and Breaker the first player (FP) and the second player (SP).
Where and under which license the archive is available is stated under ``Data and code availability'' in
the Declarations at the end of the paper.

\subsection{Requirements}\label{ssec:requirements}

The proof pipeline needs a C++17 compiler with OpenMP, \texttt{make}, and Python~3.10 or later, and no
further libraries. We used g++~13.3 and CPython~3.13 under Linux on a machine with four cores (Intel Xeon,
$2.1$\,GHz) and $15$\,GB of memory. The search needs about $600$\,MB of memory, the verifiers much less.
The optional cross-checks with the integer program additionally need the Python package \texttt{highspy}
(HiGHS~1.16 was used).

\subsection{Verifying the certificate}\label{ssec:certifying}

The following two commands, run in the unpacked archive, build the main verifier and run it; an
accepting run establishes \cref{fact:H}.
{\small
\begin{verbatim}
make bin/treecheck
bin/treecheck --table results/final/strategy.tbl \
    --leaves results/final/leaves.txt --rmax 10 --exact --paranoid \
    --threads 4 --out treecheck.json
\end{verbatim}
}
\noindent (The option \fn{--paranoid} is the default of the program and is given here only for
emphasis; it is switched off by \fn{--fast}.) The program exits with code $0$ and prints one line of
JSON on standard output, which begins
{\small
\begin{verbatim}
{"ok":true,"p_states":153135,"q_states":93164,"leaves_checked":106723,
 "max_t":10,"errors":[],"error_count":0,"complete":true,"timed_out":false,
 "root_void":false,"rmax":10,"root":"0 - -", ...
 "exact_visited_set":true,"paranoid_flows":true, ...
\end{verbatim}
}
\noindent and continues with \texttt{"states\_visited":246299}, the counters of \cref{tab:tree}, and the byte counts
and FNV checksums of the two input files. Except for the fields that record threads and timing, the output
must be identical to \fn{results/final/treecheck.json}.

On the machine described in \cref{ssec:requirements}, compilation takes about $10$ seconds and the
run less than $3$ seconds of wall-clock time with four threads (about $10$ CPU-seconds) when the cores are
not occupied by other jobs (\fn{logs/paper/artifact_test_final.log}, the test of the archive of
\cref{ssec:artifact}; the times vary by one or two seconds from run to run). The other running times of this program that are
quoted in this paper were measured on a loaded machine; they are compared in \cref{ssec:supp-timing}.

The second, independent verifier is run by
{\small
\begin{verbatim}
python3 verify/spotcheck.py --selftest
python3 verify/spotcheck.py --table results/final/strategy.tbl \
    --leaves results/final/leaves.txt --rmax 10 --samples 0 --seed 20261112 \
    --procs 4 --full --all-leaves --out spotcheck_full.json
\end{verbatim}
}
\noindent The first command (about $10$ seconds) tests the geometry, the canonical form, the leaf test and $\LOST$ against
brute force and ends with \texttt{selftest OK}. The second expands the whole tree and recomputes every leaf
(about $6$ CPU-minutes). Its last line of output must report \texttt{ok=True}, \texttt{errors=0},
\texttt{p\_expanded=46412}, \texttt{q\_checked=93164}, \texttt{void=151} and
\texttt{leaves\_recomputed=106723}, together with the numbers of children of each type given in
\cref{tab:verifiers}; the line is reproduced in full in \cref{ssec:supp-timing}. All counters in its output
file must agree with those of \fn{results/final/spotcheck_full.json}.

The command \texttt{sha256sum -c MANIFEST.sha256} checks all files of the archive. The SHA-256
checksums of the two certificate files are given in \cref{ssec:intro-artifact}.

\subsection{Reproducing everything}\label{ssec:reproduce}

The script \fn{scripts/reproduce.sh} recomputes the whole pipeline from the sources into
\fn{results/repro/}, with logs in \fn{logs/repro/}. It builds the eleven programs and runs their
self-tests, recomputes the minimum number of stones for a pairing and the leaf table of busy sets, runs
the search with its self-check and compares the regenerated certificate byte for byte with the one
distributed, runs both verifiers on the regenerated certificate, compares the leaf test of the library
with that of \fn{treecheck} on all leaves and on $20{,}000$ perturbed instances, repeats the exact
enumeration of $\CC$ on the leaves with $\theta\le 3$, and plays simulated games. It prints \texttt{PASS}
or \texttt{FAIL} for each step and ends with \texttt{ALL STEPS PASSED} if every step passed. We validated
the script end to end with $20{,}000$ simulated games and two threads: all $26$ steps passed, in about
$23$ minutes. In this run the regenerated certificate was byte-identical to the one distributed
(on determinism see \cref{ssec:search}; the script does not depend on it). The steps and their running times
(\cref{tab:reproduce}), the expected duration of a complete run and the optional checks with the integer
program are described in \cref{ssec:supp-reproduce}.

\subsection{Writing your own verifier}\label{ssec:own}

Readers who do not wish to trust our programs can verify \cref{fact:H} with a program of
their own. Everything needed is in this paper: the board and the group (\cref{sec:board}), admissibility and
the leaf condition (\cref{def:admissible,def:leaf}), the states and their children (\cref{def:states}), the
two algorithms (\cref{alg:verify,alg:leaf}), and the file formats, the canonical form and test vectors
(\cref{sec:formats}). The task is small. The two re-expansions written during the internal review have
about $700$
lines of C++ each, compute canonical forms by brute force over the $1{,}152$ maps, and run in $1.5$ to
$5$ CPU-minutes (the unpruned leaf test of the first accounts for most of the difference). For
$\Rmax=10$ the predicate $\LOST$ need not be implemented at all
(\cref{prop:lost-stones}); it is enough to check $t<10$ at every expanded P-state.

That the paper suffices was tested during the internal refereeing of its first version
(\cref{tab:verifiers}; on the nature of this refereeing see \cref{ssec:impl} and the Declarations). One
referee, who was
given only the PDF of the first version of this paper and the two certificate files, wrote a verifier of
$706$ lines of C++ (\fn{logs/paper/replication_from_paper/verify.cpp}). According to the report
(\fn{logs/paper/referee_replication.md}), the first implementation written from the text passed every
test vector and reproduced every count without a correction, except two statistics of the main verifier
that were printed in the first version without the rules behind them (the numbers of flows for pairs and
for triples, \cref{ssec:supp-mainleaf}). The program prints \texttt{ACCEPT}, visits $246{,}299$ states,
reproduces every entry of \cref{tab:tree,tab:leaf-theta}, and runs in $42$ CPU-seconds on one thread.
Its exit code is nevertheless $1$, and its log (\fn{logs/paper/replication_from_paper/run.log}) ends with
one reported mismatch: besides verifying the certificate, the program compares its own numbers of flows
with these two implementation statistics. A second referee wrote an implementation of \cref{alg:verify,alg:leaf} for an
audit of the numbers of this paper (\fn{logs/paper/referee/audit/ref_audit.cpp}; $63$ CPU-seconds), with
the same result, and a third transcribed \cref{alg:verify} from the printed pseudocode
(\fn{logs/paper/referee/b_expand.py}; no leaf test) and reproduced every number of
\cref{ssec:fmt-targets,ssec:supp-targets}. The ambiguities that the first of these referees met are listed in the report;
they have been removed in the present version
(\cref{ssec:what,ssec:algverify,sec:formats,ssec:supp-mainleaf}).

An independent verifier should reproduce the following numbers (\cref{ssec:fmt-targets}): $246{,}299$
reachable states, namely $153{,}135$ P-states ($46{,}412$ expanded, $106{,}723$ leaves) and $93{,}164$
Q-states; $143{,}498$ threat children and $28{,}275$ danger children; $151$ \VOID{} children, all of type
(ii); the counts per round of \cref{tab:tree}; and $\LEAF(B,F,t)$ true for each of the $106{,}723$ entries of
the leaf file. Agreement in all these numbers is strong evidence that two programs have expanded the
same tree; a defect that drops children shows up in them (\cref{ssec:audits}).
